%% scaffold.tex --- Planning outline for the gdmbayes MEE Application Paper
%% Compile: tectonic scaffold.tex  (or pdflatex scaffold.tex)
%%
%% This is NOT the submission draft. It is a structured scaffold showing
%% what each section needs, annotated with TODOs and open questions.
%% Full prose lives in manuscript.tex.

\documentclass[11pt,a4paper]{article}

% ---- Packages ----------------------------------------------------------------
\usepackage[T1]{fontenc}
\usepackage[utf8]{inputenc}
\usepackage[margin=2.5cm]{geometry}
\usepackage{setspace}
\usepackage{enumitem}
\usepackage{xcolor}
\usepackage{titlesec}
\usepackage{mdframed}
\usepackage{booktabs}
\usepackage[authoryear,round]{natbib}
\usepackage{hyperref}

\hypersetup{colorlinks=true, linkcolor=blue!60!black, urlcolor=blue!60!black}

% ---- Colours -----------------------------------------------------------------
\definecolor{donegreen}{HTML}{1a7a1a}
\definecolor{todored}{HTML}{c0392b}
\definecolor{questionblue}{HTML}{1a3c8c}
\definecolor{notegray}{HTML}{555555}

% ---- Annotation macros -------------------------------------------------------
\newcommand{\done}[1]{{\color{donegreen}\textbf{[DONE]} #1}}
\newcommand{\todo}[1]{{\color{todored}\textbf{[TODO]} #1}}
\newcommand{\question}[1]{{\color{questionblue}\textbf{[?]} \textit{#1}}}
\newcommand{\note}[1]{{\color{notegray}\textit{Note: #1}}}

\newenvironment{scaffold}[1]{%
  \begin{mdframed}[backgroundcolor=gray!6, linecolor=gray!30,
                   leftmargin=0pt, rightmargin=0pt,
                   innerleftmargin=8pt, innerrightmargin=8pt,
                   innertopmargin=6pt, innerbottommargin=6pt]
  \noindent\textbf{#1}\par\vspace{2pt}
  \begin{itemize}[leftmargin=1.4em, topsep=2pt, itemsep=1pt, parsep=0pt]
}{%
  \end{itemize}
  \end{mdframed}
  \vspace{4pt}
}

% ---- Title block -------------------------------------------------------------
\title{\textbf{gdmbayes} --- Application Paper Scaffold\\[0.4em]
       \normalsize\textit{Methods in Ecology \& Evolution --- Application Paper track}}
\author{Harold Horsley}
\date{Draft: \today}

% =============================================================================
\begin{document}
\maketitle
\onehalfspacing
\tableofcontents
\vspace{1em}
\noindent\rule{\linewidth}{0.4pt}

% ---- Legend ------------------------------------------------------------------
\noindent
\done{Text or section is written / data in hand.}\quad
\todo{Needs writing or analysis.}\quad
\question{Open question --- decision needed.}\quad
\note{Annotation only.}

\vspace{1em}

% =============================================================================
\section{Title and Author Block}
% =============================================================================

\begin{scaffold}{What we need}
  \item \done{Title: \textit{gdmbayes}: A Python Package for Bayesian Generalised
    Dissimilarity Modelling}
  \item \todo{Confirm co-authors (if any) and affiliations}
  \item \todo{Add institutional email / ORCID}
  \item \question{Running head --- suggest ``gdmbayes: Bayesian GDM in Python''}
  \item \note{MEE Application Papers: 3,500-word limit (excl.\ abstract,
    references, tables/figures); check compliance before submission.}
\end{scaffold}


% =============================================================================
\section{Abstract}
% =============================================================================

\begin{scaffold}{Status: mostly written --- needs Bayesian numbers}
  \item \done{Motivation: GDM widespread, no Python, no uncertainty quantification}
  \item \done{Contribution: gdmbayes, sklearn-compatible + Bayesian, PyMC}
  \item \done{Frequentist benchmark: RMSE 0.064 vs.\ 0.074 for R \texttt{gdm}}
  \item \todo{Bayesian benchmark: fill in spGDMM RMSE/MAE from CV run}
  \item \todo{Tighten to $\leq$350 words (MEE limit)}
  \item \question{Include GP spatial effect as a separate bullet point?
    Likely yes --- it is a key differentiator.}
\end{scaffold}


% =============================================================================
\section{Introduction}
% =============================================================================

\begin{scaffold}{Status: first draft written --- needs tightening}
  \item \done{GDM background --- Ferrier et al.\ 2007, I-splines, cloglog link}
  \item \done{Applications: beta diversity, conservation, climate vulnerability}
  \item \done{Gap 1 --- no Python GDM implementation exists}
  \item \done{Gap 2 --- no uncertainty quantification in any GDM tool}
  \item \done{Prior Bayesian GDM work: BBGDM (Woolley 2017), Lusiana 2023/2025,
    White et al.\ 2024 (spGDMM / NIMBLE)}
  \item \done{Contribution statement: (i) frequentist sklearn; (ii) Bayesian PyMC}
  \item \todo{Add 1-sentence road-map paragraph at end of introduction:
    ``Section~2 describes the statistical model \ldots''}
  \item \question{Cite any Python SDM packages (e.g.\ elapid) to contextualise
    the Python ecosystem gap?}
\end{scaffold}


% =============================================================================
\section{Statistical Framework}
% =============================================================================

\note{In the current draft this is folded into ``Package Description''. For MEE
the statistical model should probably be its own section so reviewers can assess
it independently of the software.}

\subsection{I-spline basis functions}

\begin{scaffold}{Status: equations written; figure missing}
  \item \done{Definition of I-splines as integral of M-splines (Ramsay 1988)}
  \item \done{Pairwise feature: $f_{k,\ell}(i,j) = |I_\ell(x_{ki}) -
    I_\ell(x_{kj})|$}
  \item \todo{FIGURE~1 --- small figure showing one predictor's spline basis
    (3 knots, degree 3): individual $I_\ell$ curves + pairwise feature illustration}
\end{scaffold}

\subsection{Frequentist GDM}

\begin{scaffold}{Status: mostly written}
  \item \done{cloglog link, IRNNLS objective}
  \item \done{Predictor importance = sum of coefficients per predictor}
  \item \todo{Note Jeffreys smoothing at $y=1$ boundary}
\end{scaffold}

\subsection{Bayesian GDM (spGDMM)}

\begin{scaffold}{Status: equations written; prior table and DAG missing}
  \item \done{Tobit (censored-normal) likelihood on $\log y$}
  \item \done{Linear predictor with $\alpha_k$ importance weights}
  \item \done{Dirichlet I-spline coefficients per predictor}
  \item \done{GP spatial random effect}
  \item \todo{Table or list of all priors: intercept, $\alpha_k$, $\sigma$,
    $\sigma_\psi$, $\ell$}
  \item \todo{Explain variance structure options: homogeneous ($\sigma^2 \sim
    \mathrm{InvGamma}$), covariate-dependent ($\sigma^2 = \exp(X_\sigma
    \beta_\sigma)$ where $X_\sigma = [1, \mathrm{poly}(d, 3)]$ uses
    orthogonal polynomials of pairwise geographic distance), and polynomial
    ($\sigma^2 = \exp(\text{cubic in } \mu)$)}
  \item \todo{Explain spatial-effect options: none / abs\_diff / squared\_diff}
  \item \question{Include a directed-acyclic graph (DAG) of the full
    hierarchical model? Likely helps MEE reviewers --- FIGURE~2}
\end{scaffold}


\subsection{Masked-holdout cross-validation}

\begin{scaffold}{Status: written}
  \item \done{Motivation: standard train/test splitting drops test sites,
    so the GP has no spatial effect at those locations.  Masked-holdout CV
    fits on \emph{all} sites but masks held-out pairs as latent variables,
    allowing the GP to sample $\psi$ at test-site locations.}
  \item \done{Pair selection: a pair $(i,j)$ is held out if \emph{either}
    endpoint belongs to the test fold, preventing test-site environmental
    information from influencing the likelihood.}
  \item \done{Likelihood split: observed pairs use a censored-Normal
    likelihood (upper-censored at $\log y = 0$, i.e.\ $y = 1$);
    held-out pairs are free Normal latent variables sharing the same
    linear predictor $\mu$ and variance $\sigma^2$.}
  \item \done{The posterior samples of the held-out latent variables are
    the predictive distribution; back-transform via
    $Z = \min(1, \exp(\log V))$ and score with RMSE, MAE, CRPS.}
  \item \todo{Explain briefly why this differs from standard pair-level
    or site-level CV (1--2 sentences).}
\end{scaffold}


% =============================================================================
\section{Package Architecture}
% =============================================================================

\begin{scaffold}{Status: mostly written --- needs architecture figure}
  \item \done{Class hierarchy: GDMPreprocessor, GDM, spGDMM, GDMModel}
  \item \done{sklearn API: fit / predict / transform}
  \item \done{Config objects: PreprocessorConfig, ModelConfig, SamplerConfig}
  \item \done{InferenceData / NetCDF serialisation via ArviZ}
  \item \todo{FIGURE~3 --- architecture diagram: class boxes with arrows for
    composition / inheritance / data flow}
  \item \todo{Code snippet: end-to-end workflow from raw data to posterior}
  \item \question{Mention SLURM/HPC scripts for large CV runs? Probably a
    supplementary note, not main text.}
\end{scaffold}


% =============================================================================
\section{Running Example: Southwest Australia Plant Diversity}
% =============================================================================

\note{Dataset: 94 sites, 856 species, 10 env.\ predictors,
$n = 4{,}371$ pairs, Bray--Curtis $y \in [0.33, 1.00]$, 138 pairs at $y=1$.}

\subsection{Dataset and cross-validation protocol}

\begin{scaffold}{Status: written}
  \item \done{Describe SW Australia dataset (Fitzpatrick 2013 / R gdm)}
  \item \done{Cite White et al.\ 2024 Table~1 as benchmark}
  \item \done{10-fold site-level CV; sites assigned to folds,
    pairs held out if \emph{either} endpoint is a test site}
  \item \done{Masked-holdout strategy for spatial models: model fitted on
    all sites, held-out pairs treated as latent variables so the GP can
    sample $\psi$ at test-site locations (see Statistical Framework)}
\end{scaffold}

\subsection{Results}

\begin{scaffold}{Status: frequentist done; Bayesian TBD; plots needed}
  \item \done{Frequentist GDM: RMSE 0.064, MAE 0.042}
  \item \todo{Run spGDMM CV for: (i) no spatial; (ii) abs\_diff;
    (iii) squared\_diff --- fill Table~\ref{tab:benchmark}}
  \item \todo{FIGURE~4 --- predicted vs.\ observed scatter, colour by CV fold,
    1:1 line, separate panels for GDM and best spGDMM}
  \item \todo{FIGURE~5 --- predictor importance posteriors (violin / HDI)
    for $\alpha_k$, NNLS importance overlaid as a point estimate}
  \item \todo{FIGURE~6 (optional) --- GP spatial effect map, posterior mean
    $\bar\psi_i$ on SW Australia site locations}
  \item \question{Report $R^2$ alongside RMSE/MAE for consistency with the
    ecology literature?}
\end{scaffold}

\begin{table}[h]
  \centering
  \caption{Predictive performance under 10-fold CV, SW Australia dataset.
    \todo{Fill spGDMM rows from MCMC CV run.}}
  \label{tab:benchmark}
  \begin{tabular}{lcc}
    \toprule
    Model & RMSE & MAE \\
    \midrule
    R \textbf{gdm} --- Ferrier \citep{white2024}              & 0.0737 & 0.0549 \\
    Best spGDMM \citep{white2024}                             & 0.0731 & 0.0545 \\
    \textbf{gdmbayes} \texttt{GDM}                           & 0.0636 & 0.0424 \\
    \textbf{gdmbayes} \texttt{spGDMM} --- no spatial effect  & [TBD]  & [TBD]  \\
    \textbf{gdmbayes} \texttt{spGDMM} --- \texttt{abs\_diff} & [TBD]  & [TBD]  \\
    \textbf{gdmbayes} \texttt{spGDMM} --- \texttt{sqr\_diff} & [TBD]  & [TBD]  \\
    \bottomrule
  \end{tabular}
\end{table}


% =============================================================================
\section{Comparison with Existing Tools}
% =============================================================================

\begin{scaffold}{Status: table written --- minor updates needed}
  \item \done{Feature comparison table: R gdm, BBGDM, Lusiana, White/NIMBLE,
    gdmbayes}
  \item \todo{Verify Lusiana 2025 (MethodsX) is a package or just a paper ---
    update ``installable'' row accordingly}
  \item \todo{Add row for ``Ocean distance support'' (gdmbayes unique feature)}
  \item \question{Should this section come before the running example?
    MEE reviewers often want the comparison early.}
\end{scaffold}


% =============================================================================
\section{Discussion}
% =============================================================================

\note{No Discussion exists yet --- this is the biggest remaining gap.}

\begin{scaffold}{Paragraph plan}
  \item \todo{Para~1 --- Restate contributions: first Python GDM, first
    packaged Bayesian GDM, sklearn integration}
  \item \todo{Para~2 --- Interpret benchmark results: why does the Python
    \texttt{GDM} beat R gdm? (IRNNLS vs.\ single-pass NNLS, Jeffreys
    smoothing, CV protocol differences)}
  \item \todo{Para~3 --- Value of Bayesian extension: UQ for predictor
    importance, censored-normal at $y=1$, GP spatial random effect}
  \item \todo{Para~4 --- Limitations: (i) MCMC cost for large $n$;
    (ii) exact GP scales as $O(n^3)$ in sites; (iii) no built-in
    multi-species or ordination link}
  \item \todo{Para~5 --- Future directions: approximate GP (HSGP),
    integration with PyMC-experimental, climate-projection workflows}
  \item \question{Mention JOSS parallel submission for open-source visibility?}
\end{scaffold}


% =============================================================================
\section{Conclusions}
% =============================================================================

\begin{scaffold}{1--2 short paragraphs}
  \item \todo{Summary: gdmbayes brings GDM into Python and adds Bayesian
    inference for the first time as a distributable package}
  \item \todo{Broader impact: UQ for biodiversity prediction, ML pipeline
    integration, reproducible config-driven workflow}
\end{scaffold}


% =============================================================================
\section{Data and Code Availability}
% =============================================================================

\begin{scaffold}{Status: mostly written}
  \item \done{GitHub URL}
  \item \done{SW Australia data from R gdm package, in \texttt{examples/data/}}
  \item \todo{Mint Zenodo DOI and add to text}
  \item \todo{PyPI release and version badge}
\end{scaffold}


% =============================================================================
\section{Figures --- Master List}
% =============================================================================

\begin{scaffold}{All figures needed for submission (target $\leq$6)}
  \item \todo{FIGURE~1 --- I-spline basis illustration: one predictor, knot
    positions, individual $I_\ell$ curves, and pairwise feature
    $|I_\ell(x_i) - I_\ell(x_j)|$. \note{1 panel; matplotlib; conceptual.}}

  \item \todo{FIGURE~2 (optional) --- DAG of the full hierarchical Bayesian
    model: plates for predictors $k$, spline knots $\ell$, pairs $(i,j)$,
    GP node over sites. \note{tikz or daft.}}

  \item \todo{FIGURE~3 --- Package architecture: class boxes
    (GDMPreprocessor, GDM, spGDMM, GDMModel) with composition /
    inheritance / data-flow arrows. \note{graphviz or draw.io.}}

  \item \todo{FIGURE~4 --- Predicted vs.\ observed dissimilarity: scatter,
    colour by CV fold, 1:1 line, panels for GDM and best spGDMM.}

  \item \todo{FIGURE~5 --- Predictor importance posteriors: violin / HDI plot
    for $\alpha_k$ (all predictors + geo), NNLS point overlaid.}

  \item \todo{FIGURE~6 (optional) --- GP spatial effect map: posterior mean
    $\bar\psi_i$ on SW Australia site locations.}
\end{scaffold}


% =============================================================================
\section{Pre-submission Checklist}
% =============================================================================

\begin{scaffold}{Items to resolve before submitting to MEE}
  \item \question{Word count: MEE Application Papers $\leq$3,500 words (main
    text excl.\ abstract, refs, tables, figures). Run \texttt{texcount}.}
  \item \question{Figure count: MEE recommends $\leq$6 figures for Application
    Papers.}
  \item \question{Confirm submit to MEE first; JOSS in parallel is an option
    (different scope, no double-submission conflict).}
  \item \todo{Add \texttt{lusiana2025} to \texttt{references.bib} and cite
    in comparison table}
  \item \todo{Run \texttt{ruff check} and \texttt{mypy} clean before
    submission (code quality claim appears in the text)}
  \item \todo{Ensure all \texttt{examples/} scripts are self-contained and
    reproducible}
  \item \todo{Mint Zenodo DOI; update GitHub repo with DOI badge}
  \item \todo{Write \texttt{CONTRIBUTING.md} and \texttt{CITATION.cff}}
  \item \todo{Add \texttt{lusiana2025} citation and verify ``installable''
    column in comparison table}
\end{scaffold}


% ---- References --------------------------------------------------------------
\bibliographystyle{plainnat}
\bibliography{references}

\end{document}
